\section{関連研究}

\subsection{ストレージ内検出とブロック I/O 研究}

ストレージベース IDS は、SAN やブロックストレージ層で侵入やファイルレベル
変化を検出する研究として始まった\cite{banikazemi2005san,allalouf2010idstor}。
これらは、ストレージ層がセキュリティ境界になり得ることを示す一方で、
ファイル意味論を失う問題も明確にした。

SSD-Insider は、SSD ファームウェア内でランサムウェアを検出し、NAND の
delayed deletion を利用して復旧する代表的な研究である\cite{baek2018ssdinsider}。
この研究は、I/O リクエストヘッダだけでもランサムウェア検出が可能であること、
また 10 秒程度の検出窓が実用的な候補になり得ることを示した。ただし、
提案器は AE ではなく軽量な特徴量と分類器に近い。

IBM Research のストレージシステム内機械学習研究は、計算ストレージ装置や
コントローラ近傍で I/O 情報からランサムウェアを検出する方向性を示す
\cite{ibmstorage2023}。また、ブロックストレージにおける汎化性の研究は、
ファイルシステム、ボリューム状態、暗号化、仮想化、benign database workload が
検出性能に強く影響することを報告している\cite{reategui2024generalizability}。
この知見は、本研究において false positive 評価と split hygiene を
中心課題に置く理由である。

DeftPunk は cloud block store における検出と snapshot/recovery を扱い、
block-store 固有の I/O 特性を利用する実用システムとして重要である
\cite{deftpunk2024}。ただし、cloud snapshot と二段階分類器を前提にするため、
本稿の組み込み AE とは deployment boundary が異なる。したがって比較対象ではなく、
block storage で false positive と復旧境界を同時に設計する先行例として扱う。
計算ストレージの標準化動向\cite{snia-csd}は、ストレージ近傍で inference を
評価する動機になるが、特定装置で volume あたり 500 KB の detector data や
多数 volume の aggregate memory/CPU が確保できる証拠にはならない。

\subsection{公開データセット}

RanSAP は、ランサムウェアと benign software のストレージアクセスパターンを
含む公開データセットであり、timestamp、LBA、size、read/write、write entropy
を含むため、初期実験に最も近い\cite{ransap,hirano2021ransap}。
RanSMAP は後続データセットとして storage access pattern に memory access pattern
を追加しており、混合実行条件で storage-only 特徴がどこまで崩れるかを調べる
診断材料になる\cite{ransmap,hirano2024ransmap}。ただし、memory access は
本研究の装置内観測境界から外れる。

benign false positive の評価には、SNIA IOTTA\cite{sniaiotta} や
UMass Storage Trace Repository\cite{umass} のブロック I/O トレースが候補になる。
NapierOne\cite{napierone} はブロック I/O トレースではないが、高 entropy な
正常ファイル、圧縮済みファイル、暗号化済みファイルを用いた controlled replay
や圧縮率校正の材料として利用できる。
WannaLaugh\cite{wannalaugh2024} のような ransomware emulator は、実 malware を
扱わずに file ordering、partial encryption、throttling などの stress condition を
増やす手段になり得る。ただし、emulator は検出器の仮定に似た traces を
生成しやすいため、主証拠ではなく robustness stress test として使う。

\subsection{エンドポイントおよびファイルシステム型検出}

UNVEIL\cite{kharraz2016unveil}、CryptoDrop\cite{scaife2016cryptodrop}、
ShieldFS\cite{continella2016shieldfs}、Redemption\cite{kharraz2017redemption}、
PayBreak\cite{kolodenker2017paybreak} などは、ランサムウェアの振る舞い検出や
復旧に関する重要な先行研究である。これらはファイルパス、プロセス、API、
ファイル内容変化、暗号鍵観測など、ブロックストレージ単体では使えない
文脈を活用する。したがって本研究の比較対象として直接利用するのではなく、
false positive、検出遅延、復旧境界、振る舞い特徴の設計に関する参考として扱う。

\subsection{時系列 AutoEncoder 異常検出}

LSTM encoder-decoder\cite{malhotra2016lstm}、spacecraft anomaly detection
\cite{hundman2018lstm}、USAD\cite{audibert2020usad} などは、
時系列を再構成し、再構成誤差を anomaly score とする方法の基礎になる。
しかし、ストレージ装置内、10 秒統計、volume あたり 500 KB detector-data 制約、
多数 volume の aggregate cost、MNN CPU 推論という条件を同時に満たす
ランサムウェア AE 研究は限定的である。本研究の novelty は
大規模深層モデルではなく、この狭い制約下での検証可能性に置く。
